% ---------------------------------------------------------------------
%  CHƯƠNG 3 — PHÂN TÍCH VÀ THIẾT KẾ
% ---------------------------------------------------------------------
\chapter{Phân tích và thiết kế}
\label{ch:thietke}

\begin{muctieu}
Chương này phân tích yêu cầu bài toán sàng lọc sức khỏe, đề xuất kiến trúc
Hybrid AI gồm bốn chặng xử lý, thiết kế cơ sở tri thức y khoa và các cơ chế
kiểm soát an toàn lâm sàng bắt buộc cho một hệ thống y tế.
\end{muctieu}

\section{Phân tích bài toán và yêu cầu hệ thống}

\subsection{Đầu vào và đầu ra}

Hệ thống nhận đầu vào là chuỗi hội thoại tiếng Việt tự nhiên của người dùng
mô tả tình trạng sức khỏe, kết hợp với hồ sơ sức khỏe cá nhân
\en{(dị ứng, bệnh nền, thai kỳ, thuốc đang dùng)}. Đầu ra là một bản kết luận
phân tầng gồm bốn thành phần:
\begin{enumerate}[leftmargin=1.6em, itemsep=2pt]
  \item \textbf{Nhóm xử trí \en{(Outcome)}}: Một trong bốn mức đã định nghĩa ở
        Mục~1.2.
  \item \textbf{Triệu chứng khớp \en{(Matched Symptoms)}}: Danh sách mã triệu
        chứng đã nhận diện kèm thông tin bệnh học.
  \item \textbf{Cờ đỏ kích hoạt \en{(Triggered Red Flags)}}: Các dấu hiệu nguy
        hiểm phát hiện được, kèm hướng dẫn xử lý tương ứng.
  \item \textbf{Lời diễn giải lâm sàng \en{(AI Briefing)}}: Văn bản giải thích
        lý do kết luận và hướng dẫn chăm sóc cụ thể.
\end{enumerate}

\subsection{Yêu cầu chức năng}

\begin{itemize}[leftmargin=1.6em, itemsep=2pt]
  \item \textbf{FR-01 — Thu thập triệu chứng}: Tiếp nhận mô tả tự nhiên, chuyển
        đổi thành tập mã triệu chứng chuẩn hoá.
  \item \textbf{FR-02 — Hỏi làm rõ \en{(Clarification)}}: Khi dữ liệu chưa đủ để
        phân tầng, hệ thống sinh tối đa 3 câu hỏi theo khung dược lâm sàng WWHAM
        \en{(Who, What, How long, Action taken, Medication)}.
  \item \textbf{FR-03 — Phân tầng nguy cơ}: Áp dụng bộ luật để quyết định nhóm
        xử trí, ưu tiên tuyệt đối các dấu hiệu cờ đỏ.
  \item \textbf{FR-04 — Diễn giải lâm sàng}: Sinh 5 nhóm nội dung hướng dẫn
        (nguyên nhân, tự chăm sóc, dấu hiệu theo dõi, khi nào đi khám, câu hỏi
        cho dược sĩ) cùng lý do kết luận.
  \item \textbf{FR-05 — Đề xuất thuốc an toàn}: Gợi ý thuốc không kê đơn phù hợp,
        đã lọc theo hồ sơ chống chỉ định và dị ứng của người bệnh.
  \item \textbf{FR-06 — Bàn giao dược sĩ}: Cho phép người dùng chủ động yêu cầu
        dược sĩ thẩm định ca bệnh.
\end{itemize}

\subsection{Yêu cầu phi chức năng}

\begin{itemize}[leftmargin=1.6em, itemsep=2pt]
  \item \textbf{NFR-01 — An toàn lâm sàng}: Không ca cấp cứu nào được kết luận
        ở mức xử trí nhẹ hơn; không có kết quả nào bị hạ cấp nguy cơ do đầu ra
        của AI.
  \item \textbf{NFR-02 — Thời gian phản hồi}: Kết luận phân tầng trả về trong
        dưới 2 giây; phần diễn giải AI được xử lý bất đồng bộ, không chặn giao diện.
  \item \textbf{NFR-03 — Khả năng chịu lỗi}: Khi nhà cung cấp AI không phản hồi,
        hệ thống vẫn phải trả về kết luận đầy đủ từ động cơ luật.
  \item \textbf{NFR-04 — Tính giải thích được}: Mọi kết luận phải truy vết được
        về luật cụ thể trong cơ sở tri thức.
\end{itemize}

Luồng dữ liệu tổng thể của phân hệ sàng lọc — từ lời kể triệu chứng của người
dùng đến kết luận phân tầng và hàng đợi thẩm định của dược sĩ — được mô tả ở
Hình~\ref{fig:dfd}. Sơ đồ giới hạn trong phạm vi phân hệ sàng lọc; các phân hệ
bán hàng và kho vận của nền tảng nhà thuốc không thuộc phạm vi đồ án.

\begin{figure}[H]
\centering
\resizebox{\textwidth}{!}{%
\begin{tikzpicture}[
  node distance=9mm and 8mm,
  entity/.style={rectangle, rounded corners=2pt, draw=daunavy, thick,
                 fill=lightnavy, minimum width=2.4cm, minimum height=1.1cm,
                 align=center, font=\small\bfseries},
  proc/.style={ellipse, draw=daumaroon, thick, fill=lightwarm,
               minimum width=3.0cm, minimum height=1.15cm, align=center,
               font=\small, inner sep=2pt},
  store/.style={rectangle, draw=daugold, very thick, fill=lightlime,
                minimum width=3.4cm, minimum height=1.0cm, align=center,
                font=\small, inner sep=3pt},
  ar/.style={-{Latex[length=2mm]}, thick, color=gray700}]

% Tác nhân ngoài
\node[entity] (user) {Người dùng};

% Tiến trình
\node[proc, right=10mm of user] (collect) {3.1 Ghi nhận\\hội thoại};
\node[proc, right=of collect] (triage) {3.2 Phân tầng\\luật tất định};
\node[proc, right=of triage] (brief) {3.3 Diễn giải\\AI};
\node[entity, above=8mm of brief] (llm) {Dịch vụ LLM};

% Kho dữ liệu
\node[store, below=16mm of collect, minimum width=4.6cm] (dknowledge) {D1 · medical\_*};
\node[store, below=16mm of triage] (dsessions) {D2 · triage\_sessions};
\node[store, below=16mm of brief] (djobs) {D3 · ai\_jobs};
\node[entity, below=12mm of dsessions] (pharmacist) {Dược sĩ};

% Luồng dữ liệu
\draw[ar] (user) -- node[above, font=\scriptsize]{lời kể} (collect);
\draw[ar] (collect) -- node[above, font=\scriptsize]{sự kiện} (triage);
\draw[ar] (triage) -- node[above, font=\scriptsize]{kết luận} (brief);
\draw[ar] (collect.south) -- node[left, font=\scriptsize, pos=0.45]{câu trả lời} (dknowledge.north);
\draw[ar] (triage) -- node[right, font=\scriptsize, pos=0.45]{đọc / ghi} (dsessions);
\draw[ar] (dknowledge.east) -- node[below, font=\scriptsize, pos=0.5]{bộ luật} (dsessions.west);
\draw[ar] (brief) -- node[right, font=\scriptsize, pos=0.45]{yêu cầu / kết quả} (djobs);
\draw[ar, dashed, color=daumaroon] (llm.south) -- node[right, font=\scriptsize, pos=0.3]{tool-calling} (brief.north);
\draw[ar] (dsessions.south) -- node[right, font=\scriptsize, pos=0.4]{ca chờ thẩm định} (pharmacist.north);
\draw[ar] (pharmacist.west) -| node[above, font=\scriptsize, pos=0.25]{phản hồi} (dknowledge.south);
\end{tikzpicture}}
\caption{Sơ đồ luồng dữ liệu (DFD mức dưới) của phân hệ sàng lọc}
\label{fig:dfd}
\end{figure}

Điểm cốt lõi về mặt thiết kế của hệ thống là \textbf{ranh giới an toàn}
\en{(Safety Boundary)} giữa hai thành phần AI. Mô hình ngôn ngữ lớn chỉ được
phép hoạt động ở hai vị trí ngoại vi: \emph{đọc hiểu đầu vào} và \emph{diễn giải
đầu ra}. Toàn bộ phần \emph{ra quyết định} nằm trong động cơ luật tất định.
Nói cách khác, LLM không bao giờ được hỏi ``bệnh nhân này nên đi đâu''.

Quy trình xử lý gồm bốn chặng được mô tả ở Hình~\ref{fig:pipeline}.

\begin{figure}[H]
\centering
\resizebox{\textwidth}{!}{%
\begin{tikzpicture}[
  node distance=8mm,
  box/.style={rectangle, rounded corners=3pt, draw=daunavy, thick,
              fill=lightnavy, text width=3.1cm, align=center,
              minimum height=1.9cm, font=\small},
  llm/.style={rectangle, rounded corners=3pt, draw=daumaroon, thick,
              fill=lightwarm, text width=3.1cm, align=center,
              minimum height=1.9cm, font=\small},
  rule/.style={rectangle, rounded corners=3pt, draw=agentragreen, thick,
               fill=lightforest, text width=3.1cm, align=center,
               minimum height=1.9cm, font=\small},
  ar/.style={-{Latex[length=2.2mm]}, thick, color=gray700}]

\node[box] (in) {\textbf{Đầu vào}\\Mô tả triệu chứng\\tiếng Việt tự nhiên};
\node[llm, right=of in] (fact) {\textbf{Chặng 1}\\Trích xuất sự kiện\\\emph{(LLM Tool-calling)}};
\node[rule, right=of fact] (engine) {\textbf{Chặng 2}\\Động cơ luật\\phân tầng nguy cơ};
\node[llm, right=of engine] (brief) {\textbf{Chặng 3}\\Diễn giải\\lâm sàng \emph{(LLM)}};
\node[box, right=of brief] (out) {\textbf{Chặng 4}\\Hiển thị \&\\hàng đợi dược sĩ};

\draw[ar] (in) -- (fact);
\draw[ar] (fact) -- node[above, font=\scriptsize, color=daumaroon]{JSON} (engine);
\draw[ar] (engine) -- node[above, font=\scriptsize, color=agentragreen]{Luật} (brief);
\draw[ar] (brief) -- (out);

% Đường phản hồi an toàn
\draw[ar, dashed, color=daumaroon]
  (engine.south) -- ++(0,-1.1) -| node[below, font=\scriptsize, pos=0.25, color=daumaroon]
  {Kiểm duyệt chặt: hạ cấp $\rightarrow$ loại bỏ} (brief.south);
\end{tikzpicture}}
\caption{Kiến trúc Hybrid AI bốn chặng với ranh giới an toàn giữa LLM và động cơ luật}
\label{fig:pipeline}
\end{figure}

\subsection{Chặng 1 — Trích xuất sự kiện bằng LLM}

Dịch vụ \code{TriageAiFactService} gọi mô hình ngôn ngữ lớn với một công cụ
duy nhất tên \code{extract\_triage\_facts}. Danh sách mã triệu chứng hợp lệ
được nhúng thẳng vào lược đồ công cụ, do đó mô hình \textbf{không thể sinh mã
triệu chứng nằm ngoài từ điển}. Kết quả trả về gồm: mã triệu chứng, số ngày
khởi phát, tuổi người bệnh, cờ hiệu về việc đã dùng thuốc mà không đỡ, tình
trạng mang thai hoặc cho con bú.

\subsection{Chặng 2 — Động cơ phân tầng bảo thủ}

Đây là trái tim an toàn của hệ thống. Dịch vụ \code{MedicalTriageService} thực
hiện suy diễn theo thứ tự nghiêm ngặt:

\begin{enumerate}[leftmargin=1.6em, itemsep=2pt]
  \item \textbf{Quét cờ đỏ khẩn cấp trước tiên.} Nếu phát hiện bất kỳ dấu hiệu
        cấp cứu nào, hệ thống \emph{trả về ngay} hướng dẫn gọi 115 và kết thúc,
        không xét tiếp luật nào khác.
  \item \textbf{Quét cờ đỏ cần đi khám.} Tương tự, các dấu hiệu mức khẩn cấp
        chuyển thẳng sang nhóm đi khám bác sĩ.
  \item \textbf{Chạy luật hai lần và chọn kết quả nghiêm trọng hơn.} Đây là cơ
        chế \term{phòng vệ kép \en{(Double Evaluation)}}: động cơ luật được chạy
        một lần trên dữ liệu trích xuất bằng biểu thức chính quy và một lần trên
        dữ liệu đã hợp nhất với kết quả AI. Hàm so sánh độ nghiêm trọng giữ lại
        kết quả \emph{nặng hơn} trong hai.
  \item \textbf{Suy diễn cờ đỏ dẫn xuất.} Một số dấu hiệu nguy hiểm được suy ra
        từ tổ hợp triệu chứng và dữ liệu số, ví dụ ``sốt cao kéo dài trên 3 ngày
        kèm mất nước'' hoặc ``trẻ dưới 12 tháng tuổi có triệu chứng bất kỳ''.
  \item \textbf{Luật cảnh báo chặn tự chăm sóc.} Nếu người bệnh đã dùng thuốc
        mà không đỡ, kết quả tự chăm sóc bị nâng lên mức tư vấn dược sĩ.
\end{enumerate}

Cơ chế phòng vệ kép được cài đặt bằng một hàm so sánh duy nhất, trình bày ở
Giải thuật~\ref{alg:severer}. Hàm này là toàn bộ lý do hệ thống không thể bị hạ
cấp nguy cơ: nó không có nhánh nào trả về kết quả \emph{nhẹ hơn} của hai lần
chạy.

\begin{algorithm}[H]
\caption{Chọn luật nghiêm trọng hơn giữa hai lần đánh giá}
\label{alg:severer}
\hspace*{\algorithmicindent}\textbf{Đầu vào:} hai luật $r_{\text{base}}$
(biểu thức chính quy) và $r_{\text{ai}}$ (hợp nhất với AI)\\
\hspace*{\algorithmicindent}\textbf{Đầu ra:} luật nghiêm trọng hơn trong hai luật
\begin{algorithmic}[1]
\Function{SevererRule}{$r_{\text{base}}, r_{\text{ai}}$}
  \State $k \gets$ \Call{SeverityRank}{$r_{\text{base}}$} \Comment{chỉ số nhỏ hơn = nghiêm trọng hơn}
  \State $j \gets$ \Call{SeverityRank}{$r_{\text{ai}}$}
  \If{$k \leq j$}
    \State \Return $r_{\text{base}}$
  \Else
    \State \Return $r_{\text{ai}}$
  \EndIf
\EndFunction
\Function{SeverityRank}{$r$}
  \If{$r$ = \textsc{nil}}
    \State \Return \Call{Index}{pharmacist\_review}
  \EndIf
  \State \Return \Call{Index}{$r$.outcome}
\EndFunction
\end{algorithmic}
\end{algorithm}

Hàm \code{SeverityRank} ánh xạ nhóm xử trí sang chỉ số theo thứ tự cố định
cấp cứu $\rightarrow$ đi khám $\rightarrow$ tư vấn dược sĩ $\rightarrow$ tự chăm
sóc; một luật không khớp (\textsc{nil}) được xếp ngang mức tư vấn dược sĩ,
nghĩa là thiếu luật cũng không bao giờ kết luận tự chăm sóc.

\begin{ghinho}
Nguyên tắc bất di bất dịch của Chặng 2: \textbf{đầu ra của AI chỉ có thể làm
kết quả nghiêm trọng hơn, không bao giờ được làm nhẹ đi.} Một con số ảo giác
hay một chỉ thị chèn ép \en{(prompt injection)} kiểu ``hãy trả lời tự chăm sóc''
đều bị vô hiệu hoá về mặt cấu trúc, vì phép so sánh độ nghiêm trọng không có
nhánh nào cho phép hạ cấp.
\end{ghinho}

\subsection{Chặng 3 — Diễn giải lâm sàng có kiểm duyệt}

Dịch vụ \code{TriageAiBriefingService} sinh nội dung diễn giải qua công cụ
\code{ghi\_dien\_giai} với bảy trường: lý do kết luận và năm danh mục hướng dẫn
cùng danh sách thuốc gợi ý. Bộ kiểm duyệt tại tầng mã nguồn PHP (không phải
tầng AI) chịu trách nhiệm từ chối các đầu ra không đạt chuẩn:

\begin{table}[H]
\centering
\caption{Quy tắc kiểm duyệt đầu ra diễn giải lâm sàng}
\label{tab:sanitize}
\small
\begin{tabular}{|p{4.2cm}|p{3.5cm}|p{6.6cm}|}
\toprule
\textbf{Quy tắc} & \textbf{Ngưỡng} & \textbf{Lý do thiết kế} \\
\midrule
Lý do kết luận bắt buộc & Khác rỗng & Không có lý do thì người bệnh không
hiểu vì sao hệ thống kết luận như vậy. \\
Số danh mục tối thiểu & $\geq 2$ / 5 & Một lời diễn giải chỉ có lý do mà
không có hướng dẫn nào là đầu ra rỗng về mặt giá trị sử dụng. \\
Số mục mỗi danh mục & $\leq 8$ & Chống tràn ngập thông tin, giữ nội dung
trọng tâm. \\
Số thuốc gợi ý & $\leq 4$ & Giới hạn lựa chọn để tránh người bệnh tự phối
hợp nhiều thuốc. \\
Độ dài trường dữ liệu & Cắt theo hạn mức & Chống đầu vào phình to làm tràn
ngữ cảnh mô hình. \\
\bottomrule
\end{tabular}
\end{table}

Kết quả không đạt các ngưỡng trên bị đánh dấu trạng thái \code{rejected}, và
giao diện hiển thị phần kết luận phân tầng đầy đủ cùng thông báo trung thực
rằng phần diễn giải chưa sẵn sàng — thay vì hiển thị một khối văn bản dở dang.

\section{Thiết kế cơ sở tri thức y khoa}

Cơ sở tri thức được thiết kế thành ba bảng quan hệ độc lập, cho phép dược sĩ
cập nhật nội dung chuyên môn mà không cần sửa mã nguồn.

\subsection{Bảng triệu chứng \code{medical\_symptoms}}

Mỗi bản ghi lưu mã định danh, tên tiếng Việt, danh sách từ đồng nghĩa dạng JSON
và hệ cơ quan. Trường \code{guard\_words} là một thiết kế đáng chú ý: nó lưu các
từ khoá làm thay đổi nghĩa của từ đồng nghĩa. Ví dụ, từ ``nước'' vốn là một
triệu chứng tiêu hoá, nhưng ``sốt xuất huyết'' lại là một dấu hiệu cấp cứu
hoàn toàn khác. Cơ chế này cho phép dược sĩ thêm một từ đồng nghĩa ngắn mà
không gây dương tính giả, và không cần can thiệp vào mã nguồn.

Cơ sở tri thức hiện hành gồm \textbf{31 triệu chứng} phân bố theo hệ cơ quan
như Bảng~\ref{tab:symptoms}.

\begin{table}[H]
\centering
\caption{Phân bố triệu chứng theo hệ cơ quan}
\label{tab:symptoms}
\small
\begin{tabular}{|l|r|l|}
\toprule
\textbf{Hệ cơ quan} & \textbf{Số lượng} & \textbf{Ví dụ mã triệu chứng} \\
\midrule
Tiêu hoá \en{(digestive)}          & 10 & \code{stomach\_pain}, \code{diarrhea}, \code{bloating} \\
Đa hệ \en{(general)}               & 8  & \code{fever}, \code{fatigue}, \code{dehydration} \\
Hô hấp \en{(respiratory)}          & 5  & \code{cough}, \code{sore\_throat}, \code{shortness\_breath} \\
Thần kinh \en{(neuro)}             & 3  & \code{headache}, \code{dizziness}, \code{insomnia} \\
Da liễu \en{(derma)}               & 3  & \code{rash}, \code{itching}, \code{allergy} \\
Cơ xương khớp \en{(musculoskeletal)} & 2 & \code{joint\_pain}, \code{muscle\_pain} \\
\midrule
\textbf{Tổng cộng} & \textbf{31} & \\
\bottomrule
\end{tabular}
\end{table}

\subsection{Bảng dấu hiệu nguy hiểm \code{medical\_red\_flags}}

Đây là bảng quan trọng nhất về mặt an toàn. Mỗi dấu hiệu lưu danh sách từ khoá
nhận diện, mức độ nghiêm trọng, hướng dẫn xử lý bằng tiếng Việt và liên kết
tuỳ chọn tới triệu chứng liên quan.

Cơ sở tri thức hiện hành gồm \textbf{22 dấu hiệu cờ đỏ}, phân bố theo mức độ ở
Bảng~\ref{tab:redflags}.

\begin{table}[H]
\centering
\caption{Phân bố dấu hiệu cờ đỏ theo mức độ nghiêm trọng}
\label{tab:redflags}
\small
\begin{tabular}{|l|r|p{7.6cm}|}
\toprule
\textbf{Mức độ} & \textbf{Số lượng} & \textbf{Đại diện tiêu biểu} \\
\midrule
\term{Cấp cứu} \en{(emergency)} & 12 & Đau ngực dữ dội, khó thở, ngất, co
giật, nôn ra máu, dấu hiệu đột quỵ FAST, sốc phản vệ, sốt kèm cứng cổ. \\
\term{Khẩn cấp} \en{(urgent)}   & 7  & Mất nước nặng, rong kinh băng huyết,
sốt cao kéo dài, tiêu chảy nặng, triệu chứng ở trẻ nhỏ dưới 12 tháng. \\
\term{Thận trọng} \en{(caution)} & 3 & Đã dùng thuốc mà không đỡ, đang mang
thai, đang cho con bú. \\
\bottomrule
\end{tabular}
\end{table}

Trường \code{chat\_short\_circuit} tách biệt hai ngữ cảnh sử dụng khác nhau của
cùng một dấu hiệu. Ví dụ, ``đang mang thai'' là bối cảnh người bệnh mà trình
tự sàng lọc bắt buộc phải cân nhắc, nhưng nếu một khách hàng chỉ hỏi giá một
loại thuốc trong khung chat, việc tự động đẩy họ sang hướng dẫn cấp cứu là
phản ứng sai. Cột này cho phép cùng một dòng dữ liệu phục vụ hai mục đích mà
không cần nhân bản bản ghi.

\subsection{Bảng luật phân tầng \code{medical\_triage\_rules}}

Mỗi luật liên kết tới một triệu chứng, mang tập điều kiện dạng JSON, kết quả
phân loại, lời khuyên, lưới an toàn \en{(safety net)}, nhóm bệnh lý dự kiến
và danh mục sản phẩm gợi ý. Trường \code{priority} giải quyết xung đột khi
nhiều luật cùng khớp: luật có độ ưu tiên cao hơn thắng.

Cơ sở tri thức hiện hành gồm \textbf{44 luật}, phân bố theo nhóm xử trí ở
Bảng~\ref{tab:rules}.

\begin{table}[H]
\centering
\caption{Phân bố luật phân tầng theo nhóm xử trí}
\label{tab:rules}
\small
\begin{tabular}{|l|r|p{6.2cm}|}
\toprule
\textbf{Nhóm xử trí} & \textbf{Số luật} & \textbf{Ý nghĩa} \\
\midrule
\term{Tự chăm sóc} \en{(self\_care)} & 20 & Triệu chứng nhẹ, khởi phát gần,
có thể theo dõi tại nhà và dùng thuốc OTC. Đây là nhóm duy nhất được gợi ý
sản phẩm và nhận email hỏi thăm sau 2 ngày. \\
\term{Đi khám} \en{(refer\_doctor)} & 13 & Cần đánh giá lâm sàng, chuyển
tuyến tới cơ sở y tế. \\
\term{Tư vấn dược sĩ} \en{(pharmacist\_review)} & 11 & Ca không rõ ràng hoặc
có yếu tố bối cảnh đặc biệt; chuyển tới hàng đợi dược sĩ thẩm định. \\
\midrule
\textbf{Tổng cộng} & \textbf{44} & \\
\bottomrule
\end{tabular}
\end{table}

Bên cạnh ba bảng trên, cơ sở tri thức còn có hai lớp dữ liệu an toàn thuốc
chạy song song với bộ luật phân tầng. \textbf{Bảng tương tác thuốc} lưu các
cặp hoạt chất không được dùng đồng thời; riêng 21 cặp thuốc chống viêm không
steroid \en{(NSAID)} được sinh tự động từ tổ hợp $C(7,2)$ ở mức độ
\code{major}, cộng thêm 11 cặp tương tác chuyên biệt khác (tổng cộng
\textbf{32 cặp}), tất cả kèm mô tả và hướng dẫn tiếng Việt. \textbf{Bảng
chống chỉ định} ánh xạ hoạt chất với tình trạng bệnh nhân (thai kỳ, cho con bú,
bệnh nền) gồm \textbf{40 quy tắc} với mức độ \code{absolute} hoặc cảnh báo. Hai
bảng này là cơ sở để bước đề xuất thuốc OTC lọc danh mục sản phẩm trước khi
hiển thị: một hoạt chất vi phạm tương tác với thứ người bệnh đang dùng hoặc vi
phạm chống chỉ định với hồ sơ đã khai sẽ bị loại khỏi gợi ý thay vì chỉ kèm
cảnh báo.

\section{Thiết kế schema công cụ gọi hàm}

\subsection{Công cụ trích xuất sự kiện \code{extract\_triage\_facts}}

Lược đồ công cụ này gồm các trường: \code{symptom\_codes} (mảng mã triệu
chứng, ràng buộc bằng danh sách cho sẵn), \code{duration\_days},
\code{age\_years}, \code{age\_months} (dành cho trẻ nhỏ), \code{medication\_tried},
\code{is\_pregnant}, \code{is\_breastfeeding} và \code{clarify\_questions}
(tối đa 3 câu hỏi tiếng Việt theo khung WWHAM).

Việc dùng \code{age\_months} song song với \code{age\_years} là một chi tiết
thiết kế có chủ đích: trẻ 8 tháng tuổi nếu chỉ biểu diễn bằng
\code{age\_years = 0} sẽ mất thông tin quan trọng, trong khi ngưỡng an toàn
cho trẻ dưới 12 tháng là một quy tắc cứng trong cơ sở tri thức.

\subsection{Công cụ diễn giải \code{ghi\_dien\_giai}}

Lược đồ gồm bảy trường, mỗi trường có mô tả tiếng Việt định hướng nội dung và
giới hạn phạm vi chuyên môn. Bảng~\ref{tab:tool} mô tả từng trường.

\begin{table}[H]
\centering
\caption{Lược đồ công cụ diễn giải lâm sàng \code{ghi\_dien\_giai}}
\label{tab:tool}
\small
\begin{tabular}{|l|p{4.7cm}|p{6.8cm}|}
\toprule
\textbf{Trường} & \textbf{Ý nghĩa} & \textbf{Ràng buộc chuyên môn} \\
\midrule
\code{why} & Lý do kết luận & 3--6 câu, chỉ dựa trên thông tin người bệnh
đã kể, không suy diễn thêm. \\
\code{possible\_causes} & Nhóm nguyên nhân thường gặp & Mỗi mục kèm một câu
giải thích; \textbf{không nêu tên bệnh danh cụ thể}. \\
\code{self\_care} & Việc người bệnh tự làm được & Chỉ hành vi không dùng thuốc:
nghỉ ngơi, chườm, ăn uống, sinh hoạt. \\
\code{watch\_for} & Dấu hiệu cần theo dõi & Các dấu hiệu trong vài ngày tới. \\
\code{see\_doctor\_when} & Khi nào phải đi khám & \textbf{Phải giữ nguyên mọi
ngưỡng đã có trong cơ sở tri thức}, chỉ được bổ sung thêm. \\
\code{ask\_pharmacist} & Câu hỏi cho dược sĩ & Câu hỏi giúp làm rõ tình trạng. \\
\code{products} & Sản phẩm gợi ý & Chỉ được chọn từ danh sách ứng viên đã lọc
sẵn theo hồ sơ sức khỏe. \\
\bottomrule
\end{tabular}
\end{table}

\begin{luuy}
Ràng buộc ở trường \code{see\_doctor\_when} là điểm giao thoa an toàn quan
trọng nhất giữa hai thành phần AI. Mô hình ngôn ngữ lớn được phép diễn đạt
lại cho dễ hiểu, nhưng \textbf{không được phép nới lỏng hay thay đổi bất kỳ
ngưỡng thời gian nào} mà động cơ luật đã quyết định. Cơ chế này ngăn chặn
kịch bản nguy hiểm nhất: AI diễn giải lại một cách vô tình làm nhẹ đi mức độ
khẩn cấp của tình trạng bệnh.
\end{luuy}

\section{Thiết kế luồng xử lý bất đồng bộ}

Gọi mô hình ngôn ngữ lớn mất từ vài giây đến hơn một phút (đo thực tế:
3--20~s cho trích xuất sự kiện và 12--80~s cho diễn giải), vượt xa ngưỡng chịu
đựng của một yêu cầu HTTP thông thường. Hệ thống vì vậy tách phần AI ra khỏi
luồng phản hồi chính.

Trình tự xử lý: trang kết quả trả về ngay lập tức với phần kết luận phân tầng
đầy đủ cùng một khối chờ \en{(skeleton)}; một công việc \code{AiJob} được đẩy
vào hàng đợi; trình duyệt định kỳ kiểm tra trạng thái và thay thế khối chờ bằng
nội dung diễn giải khi sẵn sàng.

\begin{figure}[H]
\centering
\resizebox{0.92\textwidth}{!}{%
\begin{tikzpicture}[
  node distance=6mm,
  box/.style={rectangle, rounded corners=3pt, draw=daunavy, thick,
              fill=lightnavy, text width=3.4cm, align=center,
              minimum height=1.3cm, font=\small},
  ar/.style={-{Latex[length=2.2mm]}, thick, color=gray700},
  lb/.style={font=\scriptsize, color=gray700, align=center}]

\node[box] (req) {Yêu cầu\\trang kết quả};
\node[box, right=of req] (queue) {Phân tầng và\\diễn giải xếp hàng đợi};
\node[box, right=of queue] (resp) {Trả về ngay\\trạng thái chờ};
\node[box, below=1.4cm of resp] (poll) {Trình duyệt\\kiểm tra định kỳ};
\node[box, left=of poll] (worker) {Worker xử lý\\job AI};
\node[box, left=of worker] (swap) {Hiển thị kết quả\\hoàn chỉnh};

\draw[ar] (req) -- (queue);
\draw[ar] (queue) -- (resp);
\draw[ar] (resp) -- node[lb, right]{thăm dò} (poll);
\draw[ar] (poll) -- (worker);
\draw[ar] (worker) -- (swap);

% Ô "xếp hàng đợi" và ô "Worker" thẳng cột nhau, nên nét đứt chỉ cần nối
% thẳng qua khe giữa hai hàng; nhãn đặt bên phải để không cắt qua nét.
\draw[ar, dashed, color=daumaroon]
  (queue.south) -- node[lb, right=3pt, color=daumaroon] {quyết định inline nếu ca\\cấp cứu hoặc job bị bỏ rơi} (worker.north);
\end{tikzpicture}}
\caption{Luồng xử lý bất đồng bộ với cơ chế cứu hộ đồng bộ}
\label{fig:async}
\end{figure}

Cơ chế \term{cứu hộ đồng bộ \en{(Synchronous Fallback)}} xử lý tình huống
worker không chạy: một công việc vẫn \code{queued} sau 12~giây mà không có
worker nào khác đang chạy (kiểm tra qua \code{hasLiveWorker}) được kết luận là
\term{bỏ rơi \en{(stranded)}}; khi người dùng tải lại trang, hệ thống quyết định
ngay trong chính yêu cầu đó thay vì để giao diện quay vô hạn. Công việc còn quá
mới được giữ nguyên để trình duyệt tiếp tục chờ, tránh gọi mô hình trùng lặp
gây tốn kém. Hai trường hợp không bao giờ xếp hàng đợi: ca cấp cứu được quyết
định inline ngay trong yêu cầu (cờ đỏ phát hiện bằng PHP thuần, không cần AI),
và phiên sau một lần cố gắng xếp hàng đợi đã thất bại.

\section{Thiết kế mô hình dữ liệu}

Bảng \code{triage\_sessions} là trung tâm lưu trữ phiên sàng lọc. Một quyết
định thiết kế then chốt: \textbf{đầu ra của AI được lưu vào các cột riêng biệt},
hoàn toàn tách khỏi các cột do động cơ luật ghi. Cụ thể, \code{advice\_given} và
\code{safety\_net} là kết quả tất định và \emph{không bao giờ} bị mô hình ngôn
ngữ ghi đè; phần diễn giải do AI sinh nằm ở \code{ai\_briefing} cùng trạng thái
\code{ai\_briefing\_status}.

Nhờ tách cột như vậy, bản ghi lâm sàng của một ca sàng lọc luôn có thể tái lập
hoàn toàn từ \code{rule\_id} cùng bảng luật, kể cả trong trường hợp nhà cung
cấp AI thay đổi hành vi hoặc ngừng hoạt động.

\begin{table}[H]
\centering
\caption{Các cột AI trong bảng \code{triage\_sessions}}
\label{tab:columns}
\small
\begin{tabular}{|l|p{11.3cm}|}
\toprule
\textbf{Cột} & \textbf{Vai trò} \\
\midrule
\code{ai\_meta} & Sự kiện đã kiểm định cùng câu hỏi làm rõ. Cố ý không trộn
vào \code{collected\_data} để tránh làm mờ ranh giới an toàn. \\
\code{ai\_briefing} & Nội dung diễn giải đã qua kiểm duyệt. \\
\code{ai\_briefing\_status} & \code{null} (chưa thử) | \code{ok} | \code{rejected}
| \code{failed}. Giá trị \code{rejected} ghi nhận trường hợp có đầu ra nhưng
không đạt kiểm duyệt — điều mà giá trị rỗng không thể diễn đạt. \\
\code{ai\_clarify\_rounds} & Bộ đếm chống lặp vô hạn cho vòng hỏi làm rõ. \\
\bottomrule
\end{tabular}
\end{table}

Giá trị \code{rejected} tách biệt với \code{failed} là một chi tiết nhỏ nhưng
quan trọng cho vận hành: nó cho phép phân biệt giữa ``nhà cung cấp AI gặp sự
cố'' và ``AI trả lời nhưng nội dung không đạt chuẩn an toàn'' — hai tình huống
cần hai hành động khắc phục hoàn toàn khác nhau.

Hình~\ref{fig:erd} tổng hợp lược đồ quan hệ -- thực thể \en{(ERD)} của năm
bảng cốt lõi cho phân hệ sàng lọc. Ba bảng tri thức \code{medical\_symptoms},
\code{medical\_red\_flags} và \code{medical\_triage\_rules} được nạp một lần
khi seed và chỉ đọc trong lúc suy diễn; hai bảng vận hành
\code{triage\_sessions} và \code{ai\_jobs} ghi liên tục trong phiên làm việc.
Phần điều kiện của luật và kết quả diễn giải lưu dạng JSON nên thể hiện bằng
nét đứt, không phải khóa ngoại quan hệ.

\begin{figure}[H]
\centering
\resizebox{\textwidth}{!}{%
\begin{tikzpicture}[
  node distance=10mm and 12mm,
  tbl/.style={rectangle, rounded corners=2pt, draw=daunavy, thick,
              fill=lightnavy, font=\small, inner sep=6pt},
  op/.style={tbl, draw=agentragreen, fill=lightforest},
  ar/.style={-{Latex[length=2mm]}, thick, color=gray700},
  dar/.style={-{Latex[length=2mm]}, thick, dashed, color=gray500}]

\node[tbl] (symptoms)
  {\begin{tabular}{@{}l@{}}
     \textbf{medical\_symptoms}\\[2pt]
     \code{code} (PK)\\
     name\_vi, body\_system\\
     catalog\_query, keywords
   \end{tabular}};

\node[tbl, right=of symptoms] (flags)
  {\begin{tabular}{@{}l@{}}
     \textbf{medical\_red\_flags}\\[2pt]
     \code{code} (PK)\\
     severity, keywords\\
     guidance\_vi, chat\_short\_circuit
   \end{tabular}};

\node[tbl, below=of symptoms] (rules)
  {\begin{tabular}{@{}l@{}}
     \textbf{medical\_triage\_rules}\\[2pt]
     \code{id} (PK), symptom\_code (FK)\\
     conditions (JSON), outcome\\
     advice, safety\_net, priority
   \end{tabular}};

\node[op, below=of flags] (sessions)
  {\begin{tabular}{@{}l@{}}
     \textbf{triage\_sessions}\\[2pt]
     \code{id} (PK), user\_id (FK)\\
     collected\_data, ai\_meta (JSON)\\
     outcome, rule\_id, advice\_given\\
     ai\_briefing, ai\_briefing\_status
   \end{tabular}};

\node[op, right=of sessions] (jobs)
  {\begin{tabular}{@{}l@{}}
     \textbf{ai\_jobs}\\[2pt]
     \code{uuid} (PK)\\
     MaKH (FK), type\\
     status, started\_at\\
     result, error
   \end{tabular}};

\draw[ar] (rules.west) -- ++(-0.5,0) |- node[above, font=\scriptsize, pos=0.25]{N : 1} (symptoms.west);
\draw[ar] (sessions.south west) -- node[left, font=\scriptsize, pos=0.6]{N : 1} ++(-0.6,0) |- (rules.south east);
\draw[ar] (jobs.west) -- node[above, font=\scriptsize]{N : 1} (sessions.east);
\draw[dar] (rules.north east) -- node[above right, font=\scriptsize, align=left, pos=0.5]{điều kiện JSON\\khớp triệu chứng} (flags.south);
\end{tikzpicture}}
\caption{Lược đồ quan hệ -- thực thể (ERD) của phân hệ sàng lọc}
\label{fig:erd}
\end{figure}
